% R17 - PROPOSED replacement of the original short-introduction
% related-literature positioning paragraph only, after author review.
% This standalone snippet is NOT included by the paper builder.
% Preserve all old scientific conclusions, equations, original examples
% and appendix results if the author chooses to integrate it.
%
\paragraph{Contribution relative to the closest prior work.}
Our research question is not whether private compatibility,
fixed-cost multi-standard production, government standards unions,
or private substitutes for public harmonization can arise:
all are familiar mechanisms.  In particular,
\citet{gandalShy2001} derive governments' incentives to
recognize foreign standards and form exclusionary standardization
unions in an explicitly price-competitive three-country setting.
\citet{takaradaEtAl2020} analyze a regional-versus-multilateral
welfare and coalition-stability comparison for mandatory
standards addressing negative consumption externalities, while
\citet{klimenko2009} studies trade agreements over public policies
toward technical compatibility in network markets.
\citet{buccellaEtAl2023} and \citet{buccellaEtAl2026}
analyze strategic privately chosen compatibility, the latter
under both Cournot and Bertrand competition.  Private adaptation
\emph{after} public standards choice is not new either:
\citet{maggiMrazova2024} consider firms' fixed costs of
regulatory diversity and show that international agreements
can either promote harmonization or preserve diversity.

The conditional contribution lies in a different \emph{joint
equilibrium object}.  Governments have already established a
formal regional standards partition and an outsider firm can
\emph{privately} pay once to support the member standard without
joining the partition; coalition members need not support the
outsider standard in return.  Under sufficient conditions derived
from the firms' actual investment incentives and from
consumer-optimized national incidence, that selected asymmetry
erodes a regional member's internal protection advantage while
leaving its disadvantage in the outsider market largely intact.
Consequently the member's \emph{ranking} of the unchanged regional
and international institutions reverses.  We prove a
competition-conduct-neutral sufficient mechanism and verify
nonempty implementations under separately solved Cournot
and differentiated Bertrand market games.  The original
Cournot residual-cost threshold and documented failures
remain essential limits: our theorem is not a general
prediction that private interoperability always induces
multilateral harmonization.  Nor does the ranking result
by itself demonstrate actual post-investment treaty
renegotiation or broad coalition-stability portability.

% BibTeX key buccellaEtAl2026 is proposed in
% research/r17/R17_PROPOSED_ADDITIONAL_REFERENCES.bib.
% maggiMrazova2024 is already present in paper/references.bib;
% avoid adding a duplicate BibTeX entry.
